\subsection{Related Work on Fine-Grained Quantization}

A central challenge in 8-bit mixed-precision training is the control of quantization
error that arises from the inherently limited dynamic range of the FP8 formats
(E4M3/E5M2).
To address this challenge, a family of \emph{fine-grained quantization} strategies has
been developed; these methods assign independent scaling factors to disjoint subsets of
the data so that the local value-to-bit mapping can be adjusted to minimize
quantization loss.
The principal schemes are described below.

\paragraph{Per-Token Quantization (Activations).}
In DiT-based and large language models, activations exhibit substantial magnitude
variation across the token dimension and are frequently affected by outlier values.
When a single global per-tensor scaling factor is applied to an entire activation
tensor, the outliers force the scale to be large, which compresses the majority of
well-behaved values into a narrow and densely populated region of the FP8 grid and
incurs severe quantization error.
Per-token quantization mitigates this effect by assigning an independent scaling factor
to each input token or tile, which ensures that the numerical range of every token is
adequately covered by the FP8 format~\cite{micikevicius2022fp8}.
Per-token scaling, however, introduces non-negligible computational overhead: prior to
each quantization step, a global reduction along the token dimension is required to
compute $A_{\max}$ for every token independently.
This reduction induces a serial data dependency with the downstream matrix
multiplication, which prevents the two operations from being overlapped and imposes
measurable latency on the critical path, particularly for long sequences or large batch
sizes in which the cost of the reduction grows proportionally.

\paragraph{Per-Channel and Per-Block Quantization (Weights).}
Model weights exhibit substantially different statistical distributions across output
channels, with certain channels carrying considerably larger or smaller magnitudes than
others.
Per-channel quantization assigns a dedicated scale to each output channel, which
captures channel-level distributional characteristics accurately and substantially
reduces inter-channel quantization error.
Per-block quantization refines this approach further by applying independent scales to
fixed-size rectangular tiles (e.g., $128\!\times\!128$), thereby balancing
quantization accuracy and metadata overhead.
A block size of $128\!\times\!128$ is a common choice in practice, as it captures
intra-matrix local distributional variation while keeping the number of scale metadata
entries manageable.

\paragraph{MXFP8 Microscaling.}
Proposed by Rouhani~et~al.~\cite{rouhani2023microscaling} of Microsoft, NVIDIA, AMD,
and collaborating institutions, the Microscaling FP8 (MXFP8) format adopts a finer
quantization granularity in which every 32 consecutive elements share a single 8-bit
microscale factor.
Compared with per-channel and per-block schemes, the finer granularity of MXFP8
captures local numerical fluctuations more precisely and exhibits stronger
quantization robustness in the presence of extreme outlier distributions.
MXFP8 has been incorporated into the Open Compute Project (OCP) specification and is
supported by major hardware vendors.
These advantages, however, are obtained at the cost of denser scale-metadata storage,
more frequent memory accesses for loading and broadcasting microscales, and increased
hardware implementation complexity.

\paragraph{Precision--Overhead Trade-off.}
All fine-grained quantization strategies improve numerical accuracy at the cost of
increased overhead in scale-metadata management.
Per-token, per-channel, and per-block schemes require the system to load, broadcast,
and store a substantial volume of scale values at every training or inference
iteration.
These operations consume valuable memory bandwidth, increase control-logic complexity,
and partially offset the throughput gains that 8-bit floating-point computation is
intended to deliver.
From a systems-engineering perspective, finer quantization granularity implies more
fragmented metadata access patterns and higher hardware implementation cost.

This system-level overhead constitutes the bottleneck that HiF8 is designed to
circumvent.
Owing to its wide dynamic range, which spans 38 exponent values
(Section~\ref{sec:hif8}), HiF8 can adopt coarse per-tensor quantization while still
attaining training accuracy comparable to the BF16 baseline, as demonstrated in our
OSP-Next experiments, thereby simplifying the quantization implementation and
providing tangible engineering benefits.
